\documentclass[11pt,a4paper]{article}
\usepackage[utf8]{inputenc}
\usepackage[T1]{fontenc}
\usepackage{amsmath,amssymb,amsthm,mathrsfs}
\usepackage[margin=2.3cm]{geometry}
\usepackage{tikz}
\usetikzlibrary{arrows.meta,calc,patterns,positioning}
\usepackage[most]{tcolorbox}
\usepackage{enumitem}
\usepackage{xcolor}
\usepackage{fancyhdr}
\usepackage[hidelinks]{hyperref}

\definecolor{bleu}{RGB}{31,56,100}
\definecolor{bleuclair}{RGB}{47,85,151}
\definecolor{ocre}{RGB}{150,100,0}

% --- Macros mathématiques
\newcommand{\R}{\mathbb{R}}
\newcommand{\N}{\mathbb{N}}
\newcommand{\Q}{\mathbb{Q}}
\newcommand{\Rb}{\overline{\mathbb{R}}}
\newcommand{\A}{\mathscr{A}}
\newcommand{\B}{\mathscr{B}}
\newcommand{\C}{\mathscr{C}}
\newcommand{\M}{\mathscr{M}}
\newcommand{\Pp}{\mathscr{P}}
\newcommand{\Ee}{\mathscr{E}}
\newcommand{\Qq}{\mathscr{Q}}
\newcommand{\Ss}{\mathscr{S}}
\newcommand{\Oo}{\mathscr{O}}
\newcommand{\Ff}{\mathscr{F}}
\newcommand{\Ii}{\mathscr{I}}

% --- Mise en forme
\setlength{\parindent}{0pt}
\setlength{\parskip}{6pt}
\pagestyle{fancy}
\fancyhf{}
\renewcommand{\headrulewidth}{0.4pt}
\fancyhead[L]{\small\color{bleu}MAT321}
\fancyhead[R]{\small\color{bleu}Topologie générale}
\fancyfoot[C]{\thepage}

\newtcolorbox{enonce}[2]{enhanced,breakable,colback=bleuclair!6,colframe=bleuclair,
  boxrule=0pt,leftrule=3pt,arc=0pt,fonttitle=\bfseries\color{bleu},
  coltitle=bleu,title={#1~#2},attach boxed title to top left={yshift=-2mm},
  boxed title style={colback=white,colframe=white,boxrule=0pt,left=0pt,top=0pt,bottom=0pt},
  top=4pt,bottom=4pt}

\newtcolorbox{note}{enhanced,breakable,colback=yellow!8,colframe=ocre!70,
  boxrule=0pt,leftrule=2pt,arc=0pt,left=6pt,top=2pt,bottom=2pt,
  fontupper=\small\itshape,before upper={\textbf{\upshape[Note de transcription]}\ }}

\newenvironment{preuve}{\par\textit{\textbf{Preuve :}}\par\nopagebreak[4]}{\par}
\DeclareMathOperator{\sgn}{sgn}
\newcommand{\sect}[1]{\section*{#1}\addcontentsline{toc}{section}{#1}}
\newcommand{\datepage}[1]{\begin{flushright}\small\bfseries\itshape\color{bleuclair}#1\end{flushright}\vspace{-8pt}}

\newcommand{\cE}[1]{\mathrm{C}_E^{#1}}
\begin{document}

\begin{center}
  {\Huge\bfseries\color{bleu} MAT321}\\[4pt]
  {\LARGE\bfseries\color{bleuclair} Topologie générale}\\[4pt]
  {\large\itshape Notes de cours saisies à partir du manuscrit}
\end{center}
\vspace{-4pt}
\hrule
\vspace{6pt}

\begin{note}
Les encadrés jaunes signalent une précision ajoutée, une imprécision du manuscrit ou une partie absente du cahier. Le reste du texte suit fidèlement le cours, avec sa numérotation.
Les lectures douteuses ont été remplacées par la version la plus probable.
\end{note}

%=====================================================================
\section*{Chapitre 1 : Espaces topologiques}

\subsection*{0- Introduction}

On va définir une structure mathématique (espace topologique) obtenue en adjoignant à un ensemble $E$ donné, un ensemble de parties de $E$ vérifiant des propriétés « naturelles ».

En topologie, on ne fait pas de distinction entre deux objets qui peuvent se « déformer » de façon continue l'un en l'autre.
C'est pourquoi la topologie est appelée « la géométrie du caoutchouc » (\emph{Rubber sheet topology}).
De façon précise, la topologie est l'étude des propriétés d'un espace qui sont conservées par déformations continues (qu'on appellera \emph{homéomorphismes}).

On va donner un sens précis aux notions de limite, continuité, ne pas avoir de trous, être d'un seul tenant\dots

\begin{note}
Le manuscrit écrit « Ruban sheet topology » ; je rétablis « Rubber sheet topology ». Les mots « de parties » sont barrés dans « un ensemble de parties $E$ donné ».
\end{note}

%---------------------------------------------------------------------
\subsection*{I- Topologie}

Soit $E$ un ensemble non vide, $\Pp(E)$ désigne l'ensemble des parties de $E$.

\subsubsection*{1- Définitions}

\begin{enonce}{Définition}{1}
Une \emph{topologie} sur $E$ est un ensemble $\Oo$ dont les éléments sont des parties de $E$ ($\Oo\subset\Pp(E)$) satisfaisant aux trois conditions suivantes :
\begin{description}[leftmargin=1.4cm,labelwidth=1cm,style=sameline]
  \item[$O_1$ :] l'ensemble vide $\varnothing$ et $E$ appartiennent à $\Oo$ ;
  \item[$O_2$ :] l'intersection de deux éléments de $\Oo$ appartient à $\Oo$ ;
  \item[$O_3$ :] toute réunion (finie ou non) d'éléments de $\Oo$ appartient à $\Oo$.
\end{description}
\end{enonce}

En d'autres termes, $\Oo$ est une topologie sur $E$ si :
\begin{description}[leftmargin=1.4cm,labelwidth=1cm,style=sameline]
  \item[$O_1$ :] $\varnothing,E\in\Oo$ ;
  \item[$O_2$ :] si $U,V\in\Oo$, alors $U\cap V\in\Oo$ ;
  \item[$O_3$ :] si $(U_i)_{i\in I}$ est une famille d'éléments de $\Oo$ ($\forall i\in I$, $U_i\in\Oo$), alors $\displaystyle\bigcup_{i\in I}U_i\in\Oo$.
\end{description}

\begin{enonce}{Définition}{2}
Si $\Oo$ est une topologie sur $E$, la paire $(E,\Oo)$ est appelée \emph{espace topologique}.
\end{enonce}

%---------------------------------------------------------------------
\subsubsection*{2) Exemples}

\textbf{\color{bleu}Exemple 1.} $E=\{a,b,c\}$.
\begin{align*}
  \Oo_1&=\bigl\{\varnothing,E,\{c\},\{a,b\}\bigr\},\\
  \Oo_2&=\bigl\{\varnothing,E,\{b\},\{a,b\},\{b,c\}\bigr\},\\
  \Oo_3&=\bigl\{\varnothing,E,\{a\},\{b\},\{c\},\{a,b\},\{a,c\},\{b,c\}\bigr\},\\
  \Oo_4&=\bigl\{\varnothing,E\bigr\}.
\end{align*}
$\Oo_1$, $\Oo_2$, $\Oo_3$ et $\Oo_4$ sont des topologies sur $E$.

\textbf{\color{bleu}Exemple 2.} $E=\{1,2,3,4,5,6\}$.
\[
  \Oo=\bigl\{\varnothing,E,\{1\},\{6\},\{1,6\},\{1,3,6\},\{2,3,4,5,6\}\bigr\}.
\]
$\Oo$ n'est pas une topologie sur $E$ car $\{1,3,6\}\in\Oo$, $\{2,3,4,5,6\}\in\Oo$ mais
\[
  \{1,3,6\}\cap\{2,3,4,5,6\}=\{3,6\}\notin\Oo .
\]
$O_2$ n'est pas vérifiée.

\textbf{\color{bleu}Exemple 3.} $E=\{1,2,3,4,5\}$.
\[
  \Oo=\bigl\{\varnothing,E,\{1\},\{3,4\},\{1,3,5\},\{2,3,4\}\bigr\}.
\]
$\Oo$ n'est pas une topologie sur $E$ car $\{3,4\}\in\Oo$, $\{1,3,5\}\in\Oo$ cependant
\[
  \{3,4\}\cup\{1,3,5\}=\{1,3,4,5\}\notin\Oo .
\]
$O_3$ n'est pas vérifiée.

\begin{note}
Dans l'exemple 3, $O_2$ n'est pas non plus vérifiée puisque $\{3,4\}\cap\{1,3,5\}=\{3\}\notin\Oo$ ; le manuscrit ne mentionne que $O_3$.
\end{note}

\textbf{\color{bleu}Exemple 4.} $E=\N$.
$\Oo\subset\Pp(\N)$ formé de toutes les parties finies de $\N$ et de $\N$ :
\[
  \Oo=\{\N\}\cup\{A\subset\N,\ A\text{ fini}\}.
\]
$\Oo$ n'est pas une topologie sur $\N$ car $\forall n\in\N$, $\{2n+1\}\in\Oo$ cependant
\[
  \bigcup_{n\in\N}\{2n+1\}=\{1\}\cup\{3\}\cup\dots\cup\{2n+1\}\cup\dots
\]
est infinie et $\neq\N$. Donc $\displaystyle\bigcup_{n\in\N}\{2n+1\}\notin\Oo$.

\textbf{\color{bleu}Exemple 5.} $E=\N^*$.
\[
  \Oo=\{\varnothing\}\cup\{U_n,\ n\in\N^*\}\qquad\text{où}\qquad U_n=\{p\in\N,\ p\ge n\}.
\]
On peut remarquer que $(U_n)\searrow$. $\Oo$ est une topologie sur $\N^*$.

$O_1$ : $\varnothing\in\Oo$ ; $\N^*=\{p\in\N,\ p\ge1\}=U_1\in\Oo$.

$O_2$ : $U,V\in\Oo$.
\begin{itemize}
  \item Si $U=\varnothing$ ou $V=\varnothing$, $U\cap V=\varnothing\in\Oo$.
  \item Sinon il existe $p,q\in\N^*$ tels que $U=U_p$ et $V=U_q$ ; $U\cap V=U_p\cap U_q=U_{\max(p,q)}\in\Oo$.
\end{itemize}

$O_3$ : soit $(O_i)_{i\in I}$ une famille d'éléments de $\Oo$.\\
Si $\forall i\in I$, $O_i=\varnothing$, $\bigcup O_i=\varnothing\in\Oo$.\\
Sinon $\exists j\in I$ tel que $O_j\neq\varnothing$, il existe $n_j\in\N^*$ tel que $O_j=U_{n_j}$. Soit
\[
  A=\bigl\{n_j,\ j\in I\ \big|\ \varnothing\neq O_j=U_{n_j}\bigr\}.
\]
$A$ est une partie non vide de $\N^*$. Soit $n_0=\min A$ :
\[
  \bigcup_{i\in I}O_i=\bigcup_{n\in A}U_n=U_{n_0}\in\Oo .
\]

\textbf{\color{bleu}Exemple 6} (Topologie grossière). Soit $E$ un ensemble non vide. On pose $\Oo_{\mathrm{gr}}=\{\varnothing,E\}$.
$\Oo_{\mathrm{gr}}$ est une topologie sur $E$ appelée \emph{topologie grossière} sur $E$.
L'espace topologique $(E,\Oo_{\mathrm{gr}})$ est appelé \emph{espace topologique grossier}.

\textbf{\color{bleu}Exemple 7} (Topologie discrète). Soit $E$ un ensemble non vide. On pose
\[
  \Oo_d=\Pp(E)=\{A,\ A\subset E\}.
\]
$\Oo_d$ est une topologie sur $E$ appelée \emph{topologie discrète} sur $E$.
L'espace topologique $(E,\Oo_d)$ est appelé \emph{espace topologique discret}.

\textbf{\color{bleu}Exercice} (Caractérisation de la topologie discrète). Soit $(E,\Oo)$ un espace topologique vérifiant $\forall x\in E$, $\{x\}\in\Oo$.
Montrer que $\Oo$ est la topologie discrète.

\textbf{\color{bleu}Exemple 8} (Topologie cofinie). Soit $E$ un ensemble non vide. On pose
\[
  \Oo_f=\{\varnothing\}\cup\bigl\{A\subset E\ \big|\ \cE{A}\text{ fini}\bigr\}.
\]
$\Oo_f$ est une topologie sur $E$ appelée \emph{topologie cofinie} sur $E$. En effet :

$O_1$ : $\varnothing\in\Oo_f$ ; $\cE{E}=\varnothing$ fini $\Rightarrow E\in\Oo_f$.

$O_2$ : $U,V\in\Oo_f$. Si $U=\varnothing$ ou $V=\varnothing$, alors $U\cap V=\varnothing\in\Oo_f$.
Sinon $\cE{U}$ fini et $\cE{V}$ fini ; $\cE{U\cap V}=\cE{U}\cup\cE{V}$ fini car réunion de deux ensembles finis.

$O_3$ : soit $(U_i)_{i\in I}\subset\Oo_f$. Si $\forall i\in I$, $U_i=\varnothing$, $\bigcup U_i=\varnothing\in\Oo_f$.
Sinon, $\exists j\in I$ tel que $\cE{U_j}$ fini ;
\[
  \cE{\bigcup U_i}=\bigcap_{i\in I}\cE{U_i}\subset\cE{U_j}\ \text{ fini.}
\]
$\cE{\bigcup U_i}$ est fini car sous-ensemble d'une partie finie, donc $\bigcup U_i\in\Oo_f$.

%---------------------------------------------------------------------
\subsubsection*{3) Ouverts, fermés, voisinages}

$(E,\Oo)$ est un espace topologique.

\begin{enonce}{Définition}{}
Les éléments de $\Oo$ sont appelés (ensembles) \emph{ouverts} de $(E,\Oo)$ ou ouverts de $E$ pour la topologie $\Oo$.
\[
  A\subset E\ \text{ ouvert de }(E,\Oo)\iff A\in\Oo .
\]
\end{enonce}

\begin{enonce}{Proposition}{(propriétés des ouverts)}
Soit $(E,\Oo)$ un espace topologique.
\begin{enumerate}[label=\textbf{\roman*)},leftmargin=1.2cm]
  \item $E$ et $\varnothing$ sont des ouverts.
  \item Toute intersection \underline{finie} d'ouverts est un ouvert.
  \item Toute réunion d'ouverts est un ouvert.
\end{enumerate}
\end{enonce}

\textit{Preuve.}

\begin{enonce}{Définition}{(fermé)}
Soit $(E,\Oo)$ un espace topologique. Une partie $S$ de $E$ est dite \emph{fermée} si son complémentaire dans $E$ est un ouvert.
\[
  S\subset E\ \text{ fermé de }(E,\Oo)\iff\cE{S}\in\Oo .
\]
\end{enonce}

\begin{enonce}{Proposition}{(propriétés des fermés)}
Soit $(E,\Oo)$ un espace topologique.
\begin{enumerate}[label=\textbf{\roman*)},leftmargin=1.2cm]
  \item $E$ et $\varnothing$ sont des fermés.
  \item Toute intersection de fermés est un fermé.
  \item Toute réunion finie de fermés est un fermé.
\end{enumerate}
\end{enonce}

\textit{Preuve.}

\begin{note}
Dans le manuscrit les deux propositions (propriétés des ouverts, propriétés des fermés) sont suivies du mot « Preuve » sans rien d'écrit après ; je laisse les preuves vides.
\end{note}

\textbf{\color{bleu}Exercice 2.} Caractériser les parties fermées par la topologie cofinie sur un ensemble $E$.

\begin{enonce}{Définition}{(voisinage)}
Soit $(E,\Oo)$ un espace topologique. Soit $x\in E$. Une partie $V$ de $E$ est appelée \emph{voisinage} de $x$ s'il existe une partie ouverte $O$ contenant $x$ et contenue dans $V$ :
\[
  V\text{ voisinage de }x\iff\exists O\in\Oo\ \text{tel que}\ x\in O\subset V .
\]
\end{enonce}

\textbf{Remarques :}
\begin{itemize}
  \item Si $V$ est un voisinage de $x$, $x\in V$.
  \item Si $V$ est un voisinage de $x$, $V\neq\varnothing$.
\end{itemize}
Si $O\in\Oo$ et $x\in O$, alors $O$ est un voisinage de $x$. Autrement dit, tout ouvert non vide est voisinage de chacun de ses points.

\begin{enonce}{Proposition}{(caractérisation des ouverts)}
Soit $(E,\Oo)$ un espace topologique. Une partie non vide de $E$ est un ouvert si et seulement si elle est voisinage de chacun de ses points :
\[
  \varnothing\neq A\subset E\ \text{ ouvert}\iff\forall x\in A,\ \exists O_x\in\Oo\ \text{tel que}\ x\in O_x\subset A .
\]
\end{enonce}

\textit{Preuve.}
\[
  A=\bigcup_{x\in A}O_x .
\]

\begin{note}
La preuve de la caractérisation des ouverts s'arrête à la ligne $A=\bigcup_{x\in A}O_x$ en bas de page ; la suite n'est pas sur les pages reçues.
Les exercices ne sont pas tous numérotés dans le cahier (« Exercice » pour la caractérisation de la topologie discrète, « Exercice 2 » pour les fermés de la topologie cofinie).
\end{note}

\end{document}
